\documentclass[a4paper,11pt]{article}

% ================================================================
% Packages
% ================================================================
\usepackage[margin=2.54cm]{geometry}
\usepackage{fancyhdr}
\usepackage{titlesec}
\usepackage{setspace}
\usepackage{graphicx}
\usepackage[hidelinks]{hyperref}
\usepackage{enumitem}
\usepackage{caption}
\usepackage{float}
\usepackage{array}
\usepackage{booktabs}
\usepackage{tabularx}
\usepackage{multirow}
\usepackage{mathptmx}
\usepackage{xcolor}
\usepackage{verbatim}

% ================================================================
% Page and paragraph settings
% ================================================================
\geometry{margin=2.54cm}
\singlespacing
\setlength{\parindent}{2em}
\setlength{\parskip}{0.2em}

% ================================================================
% Section title formatting
% ================================================================
\titleformat{\section}
{\Large\bfseries}
{\thesection}
{1em}
{}

\titleformat{\subsection}
{\large\bfseries}
{\thesubsection}
{1em}
{}

\titleformat{\subsubsection}
{\normalsize\bfseries}
{\thesubsubsection}
{1em}
{}

\titlespacing*{\section}
{0pt}
{1.1ex plus 1ex minus .2ex}
{0.7ex plus .2ex}

\titlespacing*{\subsection}
{0pt}
{0.9ex plus 1ex minus .2ex}
{0.45ex plus .2ex}

\titlespacing*{\subsubsection}
{0pt}
{0.7ex plus 1ex minus .2ex}
{0.25ex plus .2ex}

% ================================================================
% Header and footer
% ================================================================
\pagestyle{fancy}
\fancyhf{}
\fancyhead[C]{
    FIT5032 Internet Applications Development -- Assessed Lab 5
}
\fancyfoot[C]{\thepage}
\setlength{\headheight}{14pt}

\renewcommand{\headrulewidth}{0.4pt}
\renewcommand{\footrulewidth}{0pt}

% ================================================================
% Editable student and project information
% ================================================================
% Modify the following information before submission.

\newcommand{\studentname}{Wei Zhang}
\newcommand{\studentid}{36668699}
\newcommand{\tutorialgroup}{Tuesday 2:00 pm}
\newcommand{\tutorname}{Yujie Wang}
\newcommand{\submissiondate}{24 July 2026}

% Replace this with your actual project and GitHub repository if needed.
\newcommand{\projectname}{FIT5032-Lab5-WeiZhang}
\newcommand{\githuburl}{https://github.com/weiwei0618seu/FIT5032-Lab5-WeiZhang}

% ================================================================
% Document information
% ================================================================
\title{
    Event Handling, Two-Way Data Binding, and Vue Router Navigation
}

\author{
    \renewcommand{\arraystretch}{1.35}
    \begin{tabular}
        {|>{\bfseries}p{4.2cm}|p{8.2cm}|}
        \hline
        Student Name
        & \studentname
        \\ \hline
        Student ID
        & \studentid
        \\ \hline
        Tutorial/Lab
        & \tutorialgroup
        \\ \hline
        Project Name
        & \texttt{\projectname}
        \\ \hline
    \end{tabular}
}

\date{\submissiondate}

% ================================================================
% Custom cover page
% ================================================================
\makeatletter

\newcommand{\makecover}{
\begin{titlepage}

    \linespread{1.05}\selectfont
    \centering
    \vspace*{0.3cm}

    % Put the Monash logo file in the same folder as this .tex file.
    % Suggested filename:
    % monash-university-vector-logo.png

    \IfFileExists{monash-university-vector-logo.png}{
        \includegraphics[width=0.95\textwidth]
        {monash-university-vector-logo.png}
        \par
    }{
        \fbox{
            \parbox[c][2.8cm][c]{0.90\textwidth}{
                \centering
                \textbf{Monash University Logo Placeholder}
                \\[0.3em]
                Add
                \texttt{monash-university-vector-logo.png}
                to this folder.
            }
        }
        \par
    }

    \vspace{0.8cm}

    {\LARGE
        \textsc{Monash University}
        \par
    }

    \vspace{0.3cm}

    {\Large
        \textsc{Faculty of Information Technology}
        \par
    }

    \vspace{0.8cm}

    {\Large
        \textbf{FIT5032 Internet Applications Development}
        \par
    }

    \vspace{0.3cm}

    {\Large
        \textbf{Assessed Lab 5}
        \par
    }

    \vspace{0.5cm}

    \rule{\textwidth}{1.2pt}
    \\[0.4cm]

    {\LARGE\bfseries
        \@title
        \par
    }

    \rule{\textwidth}{1.2pt}
    \\[1.5cm]

    \vfill

    {\Large
        \@author
        \par
    }

    \vspace{1.5cm}

    {\large
        \textbf{Tutor Name:}
        \tutorname
        \par
    }

    \vspace{0.5cm}

    {\large
        \@date
        \par
    }

    \vspace*{1cm}

\end{titlepage}
}

\makeatother

% ================================================================
% Screenshot helper
% ================================================================
% Usage:
%
% \labscreenshot
% {filename.png}
% {Caption text}
% {fig:label}
%
% If the image does not exist, LaTeX displays a placeholder.

\newcommand{\labscreenshot}[3]{
\begin{figure}[H]

    \centering

    \IfFileExists{#1}{
        \includegraphics[width=0.94\textwidth]{#1}
    }{
        \fbox{
            \parbox[c][6.5cm][c]{0.90\textwidth}{
                \centering

                \textbf{Screenshot Placeholder}
                \\[0.5em]

                Expected file:
                \texttt{\detokenize{#1}}
                \\[0.5em]

                Replace this box by placing the screenshot
                in the same folder as the
                \texttt{.tex} file.
            }
        }
    }

    \caption{#2}
    \label{#3}

\end{figure}
}

% ================================================================
% Large code/file evidence helper
% ================================================================
% This is useful for Task 5.2 because the required files should be
% shown on separate pages.

\newcommand{\fileevidence}[3]{
\clearpage
\section*{File Evidence: \texttt{#1}}
\addcontentsline{toc}{section}{File Evidence: #1}

\labscreenshot
{#2}
{#3}
{fig:#2}
}

% ================================================================
% Small editable note box
% ================================================================
\newcommand{\editnote}[1]{
\begin{center}
\fbox{
    \parbox{0.90\textwidth}{
        \textbf{Before submission:} #1
    }
}
\end{center}
}

% ================================================================
% Document body
% ================================================================
\begin{document}

% Cover page without a page number
\pagenumbering{gobble}
\makecover
\newpage

% Start normal page numbering
\pagenumbering{arabic}

% ================================================================
\section{Lab Overview}
% ================================================================

This report presents the evidence completed for FIT5032 Assessed
Lab 5. The lab focuses on Vue.js event handling, two-way data
binding, Vue DevTools testing, and Vue Router navigation. The
starter library application was extended with password confirmation
validation, input-based feedback for the reason field, and routing
between different views. For the distinction and high distinction
component, custom routing and secure navigation were also
implemented to restrict selected pages based on authentication
status.

All assessed implementation work was completed individually. Peer
interaction was limited to technical discussion and mutual support
during the laboratory session.

% ================================================================
\section{Project Setup}
% ================================================================

\subsection{Activity 5.1.1: Recreating the Starter Project}

The Week 5 starter project was cloned from the official
\texttt{NoMash-Library} repository using the \texttt{week-4-end}
branch. The project was opened in Visual Studio Code, dependencies
were installed, and the development server was started using:

\begin{verbatim}
npm run dev
\end{verbatim}

If the development server had a Vite-related issue, the following
command could be used to install Vite:

\begin{verbatim}
npm install vite
\end{verbatim}

\labscreenshot
{lab5_activity511_project_running.png}
{Activity 5.1.1: The starter library application running on localhost.}
{fig:project-running}

\noindent
\textbf{Brief description:}
Figure~\ref{fig:project-running} shows that the starter library
application was running successfully before the Week 5 modifications
were added.

% ================================================================
\section{eFolio Task 5.1: Event Handling and Vue DevTools}
% ================================================================

\subsection{Activity 5.3.1: Confirm Password Field}

A new \texttt{Confirm password} input field was added to the library
registration form. The field was bound to
\texttt{formData.confirmPassword} using \texttt{v-model}, and the
corresponding error state was added to the \texttt{errors} object.

\labscreenshot
{lab5_activity531_confirm_password_form.png}
{Activity 5.3.1: Modified library form with the Confirm password input field.}
{fig:confirm-password-form}

\noindent
\textbf{Brief description:}
Figure~\ref{fig:confirm-password-form} shows the updated form layout
after the confirm password field was added.

\subsection{Password Confirmation Validation}

A custom event handler named \texttt{validateConfirmPassword} was
implemented to compare the original password and the confirm password
field. The validation was triggered using the \texttt{@blur} event,
so that the error message appears after the user finishes typing and
leaves the confirm password field.

\labscreenshot
{lab5_activity531_password_mismatch_error.png}
{Activity 5.3.1: Error message displayed when the password and confirm password fields do not match.}
{fig:password-mismatch-error}

\labscreenshot
{lab5_activity531_confirm_password_code.png}
{Activity 5.3.1: Source code implementing confirm password validation with @blur.}
{fig:confirm-password-code}

\noindent
\textbf{Brief description:}
Figures~\ref{fig:password-mismatch-error} and
~\ref{fig:confirm-password-code} show that the form displays the
message \texttt{Passwords do not match.} when the two password fields
are different.

\subsection{Activity 5.3.2: Reason for Joining Input Feedback}

The reason for joining field was also extended with input feedback.
When the user enters text containing the word \texttt{friend}, a
green message is displayed below the field:

\begin{verbatim}
Great to have a friend
\end{verbatim}

This demonstrates event-driven feedback and reactive data updates in
the form.

\labscreenshot
{lab5_activity532_reason_friend_message.png}
{Activity 5.3.2: Green feedback message displayed when the reason contains the word friend.}
{fig:reason-friend-message}

\labscreenshot
{lab5_activity532_reason_friend_code.png}
{Activity 5.3.2: Source code implementing the reason field feedback message.}
{fig:reason-friend-code}

\noindent
\textbf{Brief description:}
Figures~\ref{fig:reason-friend-message} and
~\ref{fig:reason-friend-code} show the additional feedback behaviour
for the reason for joining input field.

\subsection{Vue DevTools Testing}

Vue DevTools was used to inspect the component state during testing.
The \texttt{formData} object was checked to verify that user input
was correctly reflected in the component state through
\texttt{v-model} two-way data binding.

\labscreenshot
{lab5_devtools_component_tree.png}
{Vue DevTools showing the selected library registration form component.}
{fig:devtools-component-tree}

\labscreenshot
{lab5_devtools_formdata_state.png}
{Vue DevTools showing formData values, including password and confirmPassword.}
{fig:devtools-formdata-state}

\labscreenshot
{lab5_devtools_vmodel_update.png}
{Vue DevTools showing that form input changes are reflected in the component state.}
{fig:devtools-vmodel-update}

\noindent
\textbf{Brief description:}
Figures~\ref{fig:devtools-component-tree},
~\ref{fig:devtools-formdata-state}, and
~\ref{fig:devtools-vmodel-update} show the use of Vue DevTools to
debug and test the component state.

% ================================================================
\section{Two-Way Data Binding Evidence}
% ================================================================

\subsection{Using v-model for Form Inputs}

The form uses \texttt{v-model} to connect form inputs and component
state. When the user edits a field on the page, the related value in
\texttt{formData} is updated automatically. This confirms the use of
two-way data binding in the Vue component.

\labscreenshot
{lab5_vmodel_confirm_password_code.png}
{Source code showing v-model used for the confirm password field.}
{fig:vmodel-confirm-password-code}

\labscreenshot
{lab5_vmodel_reason_code.png}
{Source code showing v-model used for the reason for joining field.}
{fig:vmodel-reason-code}

\noindent
\textbf{Brief description:}
Figures~\ref{fig:vmodel-confirm-password-code} and
~\ref{fig:vmodel-reason-code} show the use of \texttt{v-model} for
two-way data binding in the form.

% ================================================================
\section{eFolio Task 5.2: Vue Router and Navigation}
% ================================================================

\subsection{Vue Router Setup}

Vue Router was installed to convert the one-page application into a
single-page application with multiple views. The following command
was used:

\begin{verbatim}
npm install vue-router@latest
\end{verbatim}

A \texttt{src/views} folder was created to store page-level views,
and a \texttt{src/router/index.js} file was created to define the
application routes.

\labscreenshot
{lab5_router_install_terminal.png}
{Terminal output showing Vue Router installation.}
{fig:router-install-terminal}

\labscreenshot
{lab5_views_folder_structure.png}
{Project structure showing the router folder and views folder.}
{fig:views-folder-structure}

\noindent
\textbf{Brief description:}
Figures~\ref{fig:router-install-terminal} and
~\ref{fig:views-folder-structure} show the router installation and
the project structure used for routing.

\subsection{Navigation Header with router-link}

The navigation menu was updated to use \texttt{router-link} instead
of normal anchor tags. This allows users to move between views
without refreshing the page.

\labscreenshot
{lab5_header_router_link_code.png}
{Source code showing router-link used for Home, About, and Login navigation.}
{fig:header-router-link-code}

\noindent
\textbf{Brief description:}
Figure~\ref{fig:header-router-link-code} shows that the navigation
menu uses Vue Router links for single-page application navigation.

% ================================================================
% Required file evidence for Task 5.2
% Each required file is placed on a new page with a header.
% ================================================================

\fileevidence
{src/router/index.js}
{lab5_router_index_js_code.png}
{Task 5.2: Source code of src/router/index.js showing route definitions and custom routing logic.}

\fileevidence
{src/App.vue}
{lab5_app_vue_code.png}
{Task 5.2: Source code of src/App.vue showing the router-view structure.}

\fileevidence
{src/views/HomeView.vue}
{lab5_homeview_vue_code.png}
{Task 5.2: Source code of src/views/HomeView.vue containing the registration form view.}

% ================================================================
\clearpage
\section{Localhost Routing Evidence}
% ================================================================

\subsection{Home Page Route}

The home route displays the registration form view. This confirms
that the root route is correctly connected to \texttt{HomeView.vue}.

\labscreenshot
{lab5_localhost_home_page.png}
{Application running on localhost with the Home page route open.}
{fig:localhost-home-page}

\noindent
\textbf{Brief description:}
Figure~\ref{fig:localhost-home-page} shows the application running
on the Home page route.

\subsection{About Page Route}

The about route displays the library about page. This confirms that
the \texttt{/about} route is correctly connected to the About view.

\labscreenshot
{lab5_localhost_about_page.png}
{Application running on localhost with the About page route open.}
{fig:localhost-about-page}

\noindent
\textbf{Brief description:}
Figure~\ref{fig:localhost-about-page} shows the application running
on the About page route.

% ================================================================
\section{Custom Routing and Secure Navigation}
% ================================================================

\subsection{Login Page}

A login page was created to support secure navigation. The page uses
\texttt{v-model} to bind the username and password input fields to
component state.

\labscreenshot
{lab5_custom_login_page.png}
{Custom routing: Login page with username and password fields.}
{fig:custom-login-page}

\labscreenshot
{lab5_custom_login_code.png}
{Source code implementing the Login page and hardcoded credential check.}
{fig:custom-login-code}

\noindent
\textbf{Brief description:}
Figures~\ref{fig:custom-login-page} and
~\ref{fig:custom-login-code} show the login page used for custom
routing.

\subsection{Protected Route Before Login}

The About page was configured as a protected route. When an
unauthenticated user attempts to visit the protected route, the
router redirects the user to the login page or an access-denied page.

\labscreenshot
{lab5_custom_about_redirect_before_login.png}
{Custom routing: Unauthenticated user redirected when trying to access the protected About route.}
{fig:about-redirect-before-login}

\labscreenshot
{lab5_custom_access_denied_page.png}
{Custom routing: Access denied page shown to unauthenticated users.}
{fig:access-denied-page}

\noindent
\textbf{Brief description:}
Figures~\ref{fig:about-redirect-before-login} and
~\ref{fig:access-denied-page} show that unauthenticated users cannot
access protected content.

\subsection{Successful Login and Protected Page Access}

After entering the correct hardcoded credentials, the user is
authenticated and can access the protected About page.

\labscreenshot
{lab5_custom_login_success.png}
{Custom routing: Successful login using hardcoded credentials.}
{fig:custom-login-success}

\labscreenshot
{lab5_custom_about_after_login.png}
{Custom routing: Authenticated user successfully accesses the About page.}
{fig:about-after-login}

\noindent
\textbf{Brief description:}
Figures~\ref{fig:custom-login-success} and
~\ref{fig:about-after-login} show that authenticated users can access
the protected About route.

\subsection{Logout Functionality}

A logout option was added for authenticated users. After logging out,
the user is no longer able to access protected routes.

\labscreenshot
{lab5_custom_logout_button.png}
{Custom routing: Logout option displayed for authenticated users.}
{fig:custom-logout-button}

\labscreenshot
{lab5_custom_after_logout_redirect.png}
{Custom routing: User redirected after logout and protected access removed.}
{fig:after-logout-redirect}

\noindent
\textbf{Brief description:}
Figures~\ref{fig:custom-logout-button} and
~\ref{fig:after-logout-redirect} show that logout clears the
authenticated state and protects the restricted route again.

\subsection{Navigation Guard Evidence}

A navigation guard was implemented in the router to check whether a
route requires authentication. If a route is protected and the user
is not authenticated, navigation is blocked or redirected.

\labscreenshot
{lab5_custom_navigation_guard_code.png}
{Source code showing the navigation guard used for secure routing.}
{fig:navigation-guard-code}

\noindent
\textbf{Brief description:}
Figure~\ref{fig:navigation-guard-code} shows the route guard used to
control access to protected routes.

% ================================================================
\section{Evidence Checklist}
% ================================================================

\renewcommand{\arraystretch}{1.25}

\begin{table}[H]
\centering
\begin{tabularx}
{0.96\textwidth}
{|>{\bfseries}p{3.4cm}|X|>{\centering\arraybackslash}p{2.2cm}|}
\hline
Assessment Item & Evidence Included & Status \\ \hline

Task 5.1 Password Confirmation
& Screenshot showing the modified library form and the
\texttt{Passwords do not match.} error message
& Complete \\ \hline

Task 5.1 Vue DevTools
& Screenshots showing Vue DevTools component inspection and
\texttt{formData} state changes
& Complete \\ \hline

Reason Field Feedback
& Screenshots showing the green \texttt{Great to have a friend}
message and related implementation code
& Complete \\ \hline

Two-Way Binding
& Screenshots showing \texttt{v-model} usage and data updates through
Vue DevTools
& Complete \\ \hline

Router Setup
& Screenshots showing Vue Router installation, project structure, and
navigation links
& Complete \\ \hline

Required File Evidence
& Separate-page evidence for \texttt{router/index.js},
\texttt{src/App.vue}, and \texttt{src/views/HomeView.vue}
& Complete \\ \hline

Localhost Home and About Pages
& Screenshots showing the application open on the Home page and About
page
& Complete \\ \hline

Custom Routing
& Screenshots showing login, protected route redirect, authenticated
access, logout, and navigation guard code
& Complete \\ \hline

\end{tabularx}
\caption{Assessed Lab 5 evidence checklist.}
\label{tab:evidence-checklist}
\end{table}

% ================================================================
\section{Completion Summary}
% ================================================================

The required tasks for FIT5032 Assessed Lab 5 were completed
successfully. The starter library application was recreated from the
provided Week 4 end branch and extended with additional event
handling and two-way data binding features. A confirm password field
was added to the registration form, and a custom
\texttt{@blur}-based validation function was implemented to show an
error message when the two password fields do not match. The reason
for joining field was also extended to display a green feedback
message when the input contains the word \texttt{friend}. Vue
DevTools was used to inspect and test component state, demonstrating
the use of \texttt{v-model} for two-way data binding.

For the advanced component, Vue Router was installed and configured.
The application was reorganised into views, including a Home page and
an About page. The navigation menu was updated using
\texttt{router-link}, and \texttt{router-view} was used to render the
current route. Custom secure routing was also implemented using a
login page, hardcoded credentials, protected routes, and a navigation
guard. The required screenshots and code evidence have been included
in this report.

\end{document}
